\chapter{Whose words these are}

\textbf{Purpose:} Find whose language is in this corpus, and which table holds it. \textbf{Scope:} every \texttt{.\allowbreak{}oracle.\allowbreak{}tsv} in the closed corpus, which reaches this tree under \texttt{build/\allowbreak{}oracles} through \texttt{tools/\allowbreak{}maintain/\allowbreak{}private\_\allowbreak{}sync.\allowbreak{}py}

These languages belong to the people who speak them. None of this work exists without them. Everything else in this research is measured against the tables below, and the tables are their words, written down.

Each entry opens with who spoke, because that is whose language it is. A linguist wrote the paper and a person read the paper into a table, and neither of those is whose language it is. Where a paper cites a published dictionary and never says who spoke, the entry says so, and the linguist does not go in the slot.

Conditions the speakers set are recorded with them below and hold wherever this corpus is used.

Written by \texttt{tools/\allowbreak{}Salishan/\allowbreak{}hand\_\allowbreak{}extraction/\allowbreak{}pure\_\allowbreak{}corpus\_\allowbreak{}index.\allowbreak{}py} from \texttt{corpus\_\allowbreak{}script\_\allowbreak{}extraction/\allowbreak{}paper\_\allowbreak{}config.\allowbreak{}py}, the only place a speaker's name is typed.

\subsection{Lushootseed}

\begin{itemize}
\item \textbf{Martha Lamont, Northern dialect}
\item \textbf{Annie Jack, Southern dialect}
\end{itemize}

Mellesmoen\_Kye\_ICSNL61, 467 rows, \texttt{Mellesmoen\_\allowbreak{}Kye\_\allowbreak{}ICSNL61.\allowbreak{}oracle.\allowbreak{}tsv}

Both recorded by Leon Metcalf in the 1950s. The only paper here that labels every form by dialect, and the border test is scored against it.

\subsection{Lushootseed}

\begin{itemize}
\item \textbf{Susie Sampson Peter, Upper Skagit}
\item \textbf{Martha LaMont, Tulalip-Skagit}
\end{itemize}

1983\_Hilbert, 60 rows, \texttt{1983\_\allowbreak{}Hilbert.\allowbreak{}oracle.\allowbreak{}tsv}

Vi taqʷšəblu Hilbert wrote the paper; the twenty-one examples were said by her aunt Susie Sampson Peter and by Martha LaMont, recorded by Leon Metcalf between 1950 and 1958 and by Thom Hess in 1963. The record credited Hilbert as the speaker until the hand extraction found it.

\subsection{St'át'imcets}

\begin{itemize}
\item \textbf{K̓weswapáw̓ (Linda Redan), Qayqáyten}
\item \textbf{Sam Mitchell, in van Eijk and Williams 1981}
\end{itemize}

Matthewson\_Redan\_ICSNL61, 104 rows, \texttt{Matthewson\_\allowbreak{}Redan\_\allowbreak{}ICSNL61.\allowbreak{}oracle.\allowbreak{}tsv}

K̓weswapáw̓ told the story over Zoom on 31 October 2025, three minutes twenty-eight seconds, and the audio and video are held by her. Sam Mitchell is the speaker of the earlier text the paper cites.

\subsection{St'át'imcets}

\begin{itemize}
\item \textbf{Qwa7yán'ak (Carl Alexander), Nxwísten}
\end{itemize}

AlexanderDavis\_ICSNL61, 541 rows, \texttt{Alexander\allowbreak{}Davis\_\allowbreak{}ICSNL61.\allowbreak{}oracle.\allowbreak{}tsv}

Recorded at Nxwísten on 7 July 2025, just over half an hour.

\subsection{Nuxalk}

\begin{itemize}
\item \textbf{Dr. Margaret Siwallace}
\end{itemize}

22-Nater-Bella-Coola-tale-10, 285 rows, \texttt{22-\allowbreak{}Nater-\allowbreak{}Bella-\allowbreak{}Coola-\allowbreak{}tale-\allowbreak{}10.\allowbreak{}oracle.\allowbreak{}tsv}

Recorded about 1975, published 2015.

\subsection{nɬeʔkepmxcín}

\begin{itemize}
\item \textbf{wlwlmelst (Maurice Michell), Southern yutémkt dialect}
\end{itemize}

ICSNL59\_LaFontaine\_Janzen\_final, 226 rows, \texttt{ICSNL59\_\allowbreak{}La\allowbreak{}Fontaine\_\allowbreak{}Janzen\_\allowbreak{}final.\allowbreak{}oracle.\allowbreak{}tsv}

He shares his four stories freely for people connecting with the language. They came from his mother nxwelinek and his grandmother ʔústko.

\subsection{nɬeʔkepmxcín}

\begin{itemize}
\item \textbf{Kʷəɬtəzétkʷu (Bernice Garcia), c̓əɬétkʷu (Coldwater)}
\end{itemize}

ICSNL59\_Garcia\_Hannon\_Stacey\_final, 371 rows, \texttt{ICSNL59\_\allowbreak{}Garcia\_\allowbreak{}Hannon\_\allowbreak{}Stacey\_\allowbreak{}final.\allowbreak{}oracle.\allowbreak{}tsv}

She asks it be acknowledged she is a Kamloops Indian Residential School speaker re-learning her language.

\subsection{Mainland Comox (ayajuthem)}

\begin{itemize}
\item \textbf{Mary George, Sliammon}
\item \textbf{Noel George Harry}
\item \textbf{Tommy Paul}
\end{itemize}

ICSNL56\_DavisJ\_2\_final-1, 1661 rows, \texttt{ICSNL56\_\allowbreak{}Davis\allowbreak{}J\_\allowbreak{}2\_\allowbreak{}final-\allowbreak{}1.\allowbreak{}oracle.\allowbreak{}tsv}

Recorded 1969 to 1980. The oracle check applies this paper's own repair and the coverage check applies none, which is one of the two drifts named above.

\subsection{nɬeʔkepmxcín}

\begin{itemize}
\item \textbf{Bev Phillips, Lytton First Nation (ƛ̓q̓əmcín)}
\end{itemize}

HallPhillipsICSNL60, 406 rows, \texttt{Hall\allowbreak{}Phillips\allowbreak{}ICSNL60.\allowbreak{}oracle.\allowbreak{}tsv}

Her own reading of the story is in build/audio, so it is an oracle for the extraction and not only a source. The coverage check applies a cruder space closer here than the oracle check does, the other drift named above.

\subsection{Nsyilxcən}

\begin{itemize}
\item \textbf{George Lezard, Penticton Indian Reserve}
\item \textbf{Nellie Guitterez, Upper Nicola Indian Band}
\item \textbf{Kiláwnaʔ (Andrew McGinnis), Penticton Indian Reserve}
\end{itemize}

19-Lyon\_ICSNL50\_final-78, 1417 rows, \texttt{19-\allowbreak{}Lyon\_\allowbreak{}ICSNL50\_\allowbreak{}final-\allowbreak{}78.\allowbreak{}oracle.\allowbreak{}tsv}

George Lezard recorded 1966 by Randy Bouchard, transcribed by Larry Pierre 1970, updated by permission of Arnie Baptiste, his son. Nellie Guitterez recorded 1978 or 1979 by Yvonne Hébert, reprinted by permission of Lynne Jorgesen, her great-granddaughter.

\subsection{Nsyilxcən}

\begin{itemize}
\item \textbf{Lottie Lindley, Upper Nicola}
\end{itemize}

2013\_Lindley\_Lyon, 653 rows, \texttt{2013\_\allowbreak{}Lindley\_\allowbreak{}Lyon.\allowbreak{}oracle.\allowbreak{}tsv}

\subsection{Lushootseed}

\begin{itemize}
\item The paper names no speaker. Its forms are cited from a published source.
\end{itemize}

1975\_Hilbert\_Hess, 74 rows, \texttt{1975\_\allowbreak{}Hilbert\_\allowbreak{}Hess.\allowbreak{}oracle.\allowbreak{}tsv}

A 1975 typescript, scanned, and what is on disk is OCR of the scan. The paper names no speaker for its examples.

\subsection{nɬeʔkepmxcín and Secwepemctsín}

\begin{itemize}
\item \textbf{Charley Alexis Mayoos}
\item \textbf{William Celestin}
\end{itemize}

2012\_Robertson, 195 rows, \texttt{2012\_\allowbreak{}Robertson.\allowbreak{}oracle.\allowbreak{}tsv}

Their texts are written in Chinuk pipa. Texts 3 to 6 are Chinook Jargon, which is a pidgin and is not Salish, and the who column keeps those out.

\subsection{eighteen Central Salish languages}

\begin{itemize}
\item The paper names no speaker. Its forms are cited from a published source.
\end{itemize}

WolfeICSNL60, 772 rows, \texttt{Wolfe\allowbreak{}ICSNL60.\allowbreak{}oracle.\allowbreak{}tsv}

Every form is cited from a published dictionary of one of eighteen languages, so there is nobody this corpus is of. The who column carries the language instead, and the reader writes a .pure.tsv keyed by it and no flat pure file.

\subsection{Nuxalk}

\begin{itemize}
\item The paper names no speaker. Its forms are cited from a published source.
\end{itemize}

ICSNL59\_Nater\_2\_final, 671 rows, \texttt{ICSNL59\_\allowbreak{}Nater\_\allowbreak{}2\_\allowbreak{}final.\allowbreak{}oracle.\allowbreak{}tsv}

Nater's own records, from his 1990 dictionary and 1984 grammar. No speaker is named. Six entries and two tables are Heiltsuk, Oowekyala, Kwak̓wala and Haisla, which are North Wakashan and not Salish at all.

\subsection{Nsyilxcən}

\begin{itemize}
\item \textbf{ɬk̓mxnalqs (Delphine Derrickson-Armstrong), stq̓aʔtkʷɬniw̓t}
\item \textbf{c̓əskʕáknaʔ (Dave Michele), stq̓aʔtkʷɬniw̓t}
\end{itemize}

LyonICSNL60\_Inch-2, 537 rows, \texttt{Lyon\allowbreak{}ICSNL60\_\allowbreak{}Inch-\allowbreak{}2.\allowbreak{}oracle.\allowbreak{}tsv}

Elicited from both speakers. Most cells of its two tables are starred, which is a form the linguist built and the speakers rejected, and those are held out.

\subsection{Twana}

\begin{itemize}
\item The paper names no speaker. Its forms are cited from a published source.
\end{itemize}

Kim\_TwanaReduplication\_final, 358 rows, \texttt{Kim\_\allowbreak{}Twana\allowbreak{}Reduplication\_\allowbreak{}final.\allowbreak{}oracle.\allowbreak{}tsv}

Every Twana form is Drachman's, out of a 1969 dissertation the paper calls the only reliable reference in existence for this. No speaker is named. Footnote 9's four Moses-Columbian forms are the only place its extraction repeats a combining mark, and those are corrected from the page.

\subsection{Nuxalk}

\begin{itemize}
\item The paper names no speaker. Its forms are cited from a published source.
\end{itemize}

2013\_Nater, 2605 rows, \texttt{2013\_\allowbreak{}Nater.\allowbreak{}oracle.\allowbreak{}tsv}

1407 numbered entries out of Nater's own 1990 dictionary. No speaker is named.

\subsection{nɬeʔkepmxcín}

\begin{itemize}
\item \textbf{Bev Phillips}
\item \textbf{c̓úʔsinek (Marty Aspinall)}
\item \textbf{kʷaɬtèzetkʷ (Bernice Garcia)}
\end{itemize}

Hall-et-al\_-ICSNL\_61-1, 494 rows, \texttt{Hall-\allowbreak{}et-\allowbreak{}al\_\allowbreak{}-\allowbreak{}ICSNL\_\allowbreak{}61-\allowbreak{}1.\allowbreak{}oracle.\allowbreak{}tsv}

The forms are cited from Thompson and Thompson's grammar and dictionary. These three are the speakers the paper thanks, two examples are Bev Phillips reading her own story, and kʷaɬtèzetkʷ introduces herself in the acknowledgement.

\subsection{St'át'imcets}

\begin{itemize}
\item \textbf{Qwa7yán'ak (Carl Alexander), Nxwísten}
\end{itemize}

ICSNL58\_Davis\_Mellesmoen\_final, 441 rows, \texttt{ICSNL58\_\allowbreak{}Davis\_\allowbreak{}Mellesmoen\_\allowbreak{}final.\allowbreak{}oracle.\allowbreak{}tsv}

Its data has three sources: van Eijk's dictionary, Davis et al. in preparation, and elicitation with Carl Alexander. It labels forms (U) and (L) for Upper and Lower St'át'imcets, the second external dialect label in the archive.

---

20 tables, 12338 rows read by hand.

